\chapter{Цель и задачи}
\label{ch:intro}

\section*{\textbf{Цель:}} 

Ознакомиться с основами работы в Simulink. Приобрести
практические навыки работы в Simulink, изучить его возможности для
моделирования и анализа динамических систем. Построить и исследовать модель простейшей системы массового
обслуживания (СМО) типа M/M/1 с учётом вероятности отказа
обслуживающего устройства и реализации резервирования. Исследовать, как отказы в обслуживающем устройстве и применение
резервирования влияют на ключевые показатели системы:

\begin{enumerate}
    \item Среднее время ожидания заявок в очереди
    \item Длину очереди
    \item Вероятность отказа в обслуживании заявок
\end{enumerate}

\section*{Задание к лабораторной работе}

\begin{figure}[H]
    \centering
    \includegraphics[width=1.0\textwidth]{z1.png}
    \caption{Задание для лабораторной работы}
\end{figure}

\begin{figure}[H] 
    \centering
    \includegraphics[width=1.0\textwidth]{z2.png}
    \caption{Задание для лабораторной работы}
\end{figure}

\begin{figure}[H]
    \centering
    \includegraphics[width=1.0\textwidth]{z3.png}
    \caption{Задание для лабораторной работы}
\end{figure}

\begin{figure}[H]
    \centering
    \includegraphics[width=1.0\textwidth]{z4.png}
    \caption{Задание для лабораторной работы}
\end{figure}


\endinput